\documentclass[11pt]{article}

\usepackage[margin=1in]{geometry}
\usepackage{amsmath,amssymb,amsfonts}
\usepackage{booktabs}
\usepackage{graphicx}
\usepackage{enumitem}
\usepackage{xcolor}
\usepackage{hyperref}
\hypersetup{colorlinks=true, linkcolor=blue, citecolor=blue, urlcolor=blue}
\setlist[itemize]{leftmargin=1.5em, itemsep=0.2em, topsep=0.2em}
\setlist[enumerate]{leftmargin=1.7em, itemsep=0.3em, topsep=0.3em}

\numberwithin{equation}{section}
\numberwithin{table}{section}
\numberwithin{figure}{section}

\begin{document}

\begin{abstract}
Vision-language-action models (VLAs) are typically fine-tuned for a fixed number of training steps
without an explicit validation-based stopping criterion, because training loss is a poor proxy for
closed-loop task performance and online rollout evaluation is too expensive to run at every checkpoint.
This work benchmarks offline validation metrics --- Dynamic Time Warping, Optimal Transport, cosine
distance, and an architecture-appropriate likelihood or loss surrogate --- as early-stopping signals
across five LoRA fine-tuning runs of three VLA architectures: the discrete-token OpenVLA and the
flow-matching SmolVLA and $\pi_{0.5}$. Each metric is scored against closed-loop LIBERO-goal success by
rank correlation, selection regret, and bootstrap confidence interval. Trajectory-similarity metrics
correlate positively with closed-loop success in every run ($\rho=0.28$--$0.90$), reliably identify the
converged training region, and support low-regret early-stopping checkpoint selection. Likelihood- and
loss-based surrogates are reliable only for the discrete-token architecture, and are weak or uninformative
for the flow-matching policies. No catastrophic action-quality overtraining is observed in any run, and at
the sample sizes available, statistical power rather than metric identity is usually the binding
constraint on how finely individual metrics can be ranked.
\end{abstract}

\section{Introduction}
\label{sec:intro}

Motivated by the capacity of foundation models to generalize across many tasks and domains, applying
foundation models to robotic manipulation has become an active and rapidly growing area of research
\cite{vlasurvey}. A Visual Language Model (VLA) is a large-scale generative model for robotic
manipulation that perceives its environment and task from multimodal inputs --- images, proprioceptive
state, and language instructions --- and outputs actions that interact with the environment to achieve a
specified goal. Such models are designed to adapt to real-time changes in the environment and to
generalize across tasks, environments, objects, and robot embodiments.

Training an VLA is typically framed as an imitation-learning problem, a form of supervised learning in
which a policy is trained to reproduce expert demonstrations. As in any supervised setting, each training
step updates the model's parameters and thereby its performance, so training must be stopped at an
appropriate point: too many steps risk overfitting, and too few risk underfitting \cite{ying2019overview}.
The number of steps required for convergence is affected by several interacting factors ---
dataset size, learning rate, batch size, etc. Thus, a fixed training budget chosen in
advance is not guaranteed to coincide with the point of best generalization. State-of-the-art
VLAs~\cite{kim2024openvla,shukor2025smolvla,black2024pi0} are typically fine-tuned for a fixed number of
steps with no explicit validation-based stopping criterion. As a solution, early stopping is a widely used technique in supervised learning, 
which involves monitoring a validation metric throughout training and retaining the checkpoint where the metric
 is best \cite{goodfellow2016deep}. 

Introducing early stopping for VLAs therefore requires validation metrics that can quantitatively track
model performance throughout training without relying on training loss alone. The most direct ground
truth is the task success rate observed by executing the policy online, in closed loop with the
environment. Online evaluation is, however, expensive and slow, risks damaging physical hardware before a
policy reaches an acceptable performance level, and is impractical to repeat at every one of the
potentially hundreds of thousands of steps a fine-tuning run may take. Offline validation offers a
middle ground: cheaper and faster than an online rollout, yet more predictive of downstream quality than
training loss alone, because it measures the output action similarity to expert demonstrations on a different distribution with the training set following the standard evaluation
setting of imitation learning, while avoid expensive interaction with the environment \cite{il}. Several such metrics are already used, in
one form or another, by existing VLAs: (1)~cross-entropy loss measures the divergence between the
predicted and ground-truth distributions over discrete action tokens; (2)~token-wise accuracy is
the fraction of predicted action tokens that exactly match their ground-truth counterparts; and
(3)~$L_1$/$L_2$ loss measures the absolute or squared difference between the predicted and
ground-truth action. All three, however, score actions independently at each individual time step and therefore cannot capture the temporal structure of
a trajectory: they are blind to the sequential dependencies and continuity that characterize a physically
realizable robot motion. Addressing this limitation requires metrics that jointly consider the spatial and
temporal structure of a trajectory as a whole. Dynamic Time Warping (DTW)~\cite{dtw_il}, Optimal Transport
(OT)~\cite{ot_il}, and trajectory smoothness~\cite{smoothness_il} are three such metrics already used
elsewhere in imitation-learning evaluation while not for VLA validation, each capturing temporal alignment, geometric similarity, or
motion continuity over an entire trajectory rather than at a single step, and are consequently promising candidates for offline checkpoint
selection in VLA fine-tuning. 

This work benchmarks offline early-stopping metrics against the closed-loop success rate they are
intended to predict, across three representative VLA architectures --- the discrete-token
OpenVLA~\cite{kim2024openvla} and the flow-matching SmolVLA~\cite{shukor2025smolvla} and
$\pi_{0.5}$~\cite{black2025pi05} --- fine-tuned with LoRA on the LIBERO-goal benchmark
\cite{liu2023libero} under five training configurations spanning both architecture and training-set size
(Table~\ref{tab:runs}, \S\ref{sec:implementation-runs}). Of the candidate metrics surveyed above, this
study focuses on DTW, OT, and cosine distance --- which admit a natural, model-agnostic definition
applicable identically to discrete-token and flow-matching policies alike --- together with each
architecture's most appropriate likelihood or loss surrogate, such as token negative log-likelihood, $L_1$/$L_2$,
or the flow-matching validation loss. Concretely, \S\ref{sec:method} formalizes these metrics together with the
early-stopping rule, the rank-correlation and regret statistics used to score a metric against closed-loop
success, and a bootstrap procedure for quantifying how much of any apparent metric ranking is attributable
to sampling noise; \S\ref{sec:implementation} details how this protocol was instantiated for the three
architectures; and \S\ref{sec:results} reports which offline metrics track closed-loop success most
reliably, whether they can be trusted to detect training convergence and select a low-regret checkpoint,
and where --- architecturally or statistically --- that trust breaks down.

\bigskip
\noindent\textit{Sections 2 and 3 (related work and background) are omitted from this excerpt. The
remainder of this chapter presents method (\S\ref{sec:method}), implementation
(\S\ref{sec:implementation}), and results (\S\ref{sec:results}).}

\setcounter{section}{3}

\section{Method}
\label{sec:method}

\subsection{Problem Statement and Notation}
\label{sec:method-notation}

Fine-tuning a VLA policy produces a sequence of checkpoints $\theta_0,\theta_1,\dots,\theta_K$, each
inducing a policy $\pi_k=\pi_{\theta_k}$. Selecting which checkpoint to deploy is an early-stopping
problem: the only trustworthy signal of deployment quality is closed-loop task success $Q(k)$, but
measuring $Q(k)$ for every checkpoint requires either a real robot or a costly simulator rollout. This
chapter's central question is therefore
\begin{quote}
  \emph{which offline validation metric, computed cheaply from held-out demonstration data at every
  checkpoint, best predicts the closed-loop success rate $Q(k)$ that would be observed only by
  evaluating that checkpoint online?}
\end{quote}
We write $A^\star=(a^\star_0,\dots,a^\star_{H-1})$ for an expert action chunk of horizon $H$ drawn from
a held-out validation trajectory, and $A_k$ for the corresponding chunk produced by checkpoint $k$. For
a one-step (token-autoregressive) policy such as OpenVLA, $A_k$ is formed by \emph{teacher forcing}: the
policy is queried at every expert state $h^\star_{t+j}$ along the window, $a_{k,j}=\pi_k(\,\cdot\mid
h^\star_{t+j})$. For a chunking policy such as SmolVLA or $\pi_{0.5}$, $A_k=\pi_k(\cdot\mid h^\star_t)$
is produced natively as a single forward pass. Both settings measure open-loop action quality on expert
inputs; neither captures closed-loop compounding error, a scope limitation discussed in
\S\ref{sec:implementation-deviations}.

\subsection{Offline Validation Metrics}
\label{sec:method-metrics}

Four families of metric are considered, all normalized so that \emph{lower is better}. Actions are
standardized per dimension using training-split statistics before any distance is computed,
$\bar a_d=(a_d-\mu_d)/(\sigma_d+\epsilon)$, so that translation, rotation, and gripper channels
contribute comparably.

\paragraph{Dynamic Time Warping (DTW).} DTW admits nonlinear temporal alignment between two sequences,
which is appropriate when the same underlying behavior may be executed at a slightly different pace
\cite{sakoe1978dtw}. With pointwise cost $c(u,v)=\lVert u-v\rVert_2^2$ and $\Gamma$ ranging over
monotone warping paths,
\begin{equation}
  D_{\mathrm{DTW}}(X,Y) = \frac{1}{|\Gamma^\star|}\min_{\Gamma}\sum_{(r,u)\in\Gamma} c(x_r,y_u).
  \label{eq:dtw}
\end{equation}

\paragraph{Optimal Transport (OT).} Where DTW enforces a single monotone correspondence, entropic OT
compares the two sequences as empirical distributions over their elements, admitting soft many-to-many
correspondence \cite{cuturi2013sinkhorn}. With $p,q$ uniform over the sequence elements and $\Pi(p,q)$
the transport polytope,
\begin{equation}
  D^{\varepsilon}_{\mathrm{OT}}(X,Y) = \min_{\Gamma\in\Pi(p,q)} \sum_{r,u}\Gamma_{ru}\,c(x_r,y_u)
    + \varepsilon\sum_{r,u}\Gamma_{ru}\big(\log\Gamma_{ru}-1\big).
  \label{eq:ot}
\end{equation}

\paragraph{Cosine distance.} Cosine distance discards magnitude and measures only directional
agreement per timestep,
\begin{equation}
  D_{\cos}(X,Y) = 1-\frac{1}{H}\sum_{j=0}^{H-1}\frac{x_j^\top y_j}{\lVert x_j\rVert_2\lVert y_j\rVert_2+\epsilon}.
  \label{eq:cos}
\end{equation}
Because it is magnitude-blind it is used only in conjunction with, never in place of, the metrics above.

\paragraph{Likelihood and architecture-specific surrogates.} The methodologically correct quantity is
the log-likelihood the policy assigns to the \emph{expert} action under the \emph{expert} state,
$\log\pi_k(a^\star_t\mid h^\star_t)$, not the likelihood of the policy's own rollout (an overconfident
but wrong policy can assign high likelihood to its own errors). Not every architecture exposes a
tractable likelihood, so the correct surrogate is architecture-dependent: token negative log-likelihood
for a discrete-token policy, $L_1$/$L_2$ validation error for a deterministic continuous policy, and the
validation flow-matching (or diffusion) objective for a flow-matching policy. This distinction is
load-bearing for the results in \S\ref{sec:results}: OpenVLA's NLL is a genuine token log-likelihood,
whereas SmolVLA's and $\pi_{0.5}$'s ``NLL'' is the flow-matching validation loss standing in for one,
reported under the same key so the two are visually comparable but are not the same statistical object.

\subsection{Trajectory-Level Aggregation and Task-Discriminative Validity}
\label{sec:method-taskdisc}

A metric computed on a single $H$-step chunk is a high-variance, low-context sample of a trajectory that
may span hundreds of steps. Two aggregation strategies were compared before adopting one for the
checkpoint-level score $\bar m(k)$: \textbf{(A) single-chunk}, which scores one chunk in isolation, the
atomic unit an online chunking rollout actually executes; and \textbf{(B) trajectory-level}, which
represents a trajectory by its ordered sequence of non-overlapping $H$-step chunks and compares two
trajectories \emph{window-by-window at aligned positions}, averaging the per-window distance.

Neither estimator's quality can be assumed; both must first satisfy a prerequisite that is
model-independent and precedes any question about success correlation: a validation distance is
uninformative about action quality unless it can already tell that two \emph{ground-truth} trajectories
of the \emph{same} task are more similar to one another than to trajectories of a \emph{different} task.
This is measured directly. Held-out chunks are partitioned by task and source trajectory; for each metric
we draw $N_{\mathrm{intra}}$ pairs of chunks from the same task but different trajectories, and
$N_{\mathrm{inter}}$ pairs from different tasks, and compute
\begin{equation}
  \mathrm{AUROC} = \Pr\big(d_{\mathrm{inter}} > d_{\mathrm{intra}}\big),
  \qquad
  d_{\mathrm{Cohen}} = \frac{\mu_{\mathrm{inter}}-\mu_{\mathrm{intra}}}{\sigma_{\mathrm{pooled}}},
  \label{eq:taskdisc}
\end{equation}
where $\mathrm{AUROC}=0.5$ indicates chance separation (no task information in the metric) and
$\mathrm{AUROC}=1.0$ indicates perfect separation. This test is applied identically under estimators (A)
and (B); whichever is more task-discriminative on ground truth is adopted for every $\bar m(k)$ reported
in \S\ref{sec:results}.

\subsection{Early-Stopping Selection Rule}
\label{sec:method-selection}

Given a checkpoint-level score $\bar m(k)$, the induced early-stopping rule selects the \emph{earliest}
checkpoint within a small tolerance of the metric's best observed value, rather than the checkpoint that
literally minimizes $\bar m$ (which would overfit to a single noisy point):
\begin{equation}
  k_m^\star = \min\Big\{k : \bar m(k) \le \min_j \bar m(j) + \eta_m\Big\}, \qquad
  \eta_m = 0.01\,\big|\bar m(0)-\min_j \bar m(j)\big|.
  \label{eq:selection}
\end{equation}
For a metric that is higher-is-better (token accuracy is the only such case in this study) the
inequality and the arguments to $\min$/$\max$ in Eq.~\ref{eq:selection} are reversed.

\subsection{Evaluating a Metric: Rank Correlation and Selection Regret}
\label{sec:method-eval}

A candidate metric $m$ is judged against the closed-loop success curve $Q(k)$ by two complementary
statistics, computed only over the checkpoints at which both $\bar m(k)$ and $Q(k)$ are available:
\begin{equation}
  \rho_m = \mathrm{Spearman}\big(-\bar m(k),\,Q(k)\big),
  \qquad
  \mathrm{Regret}(m) = \max_k Q(k) - Q\big(k_m^\star\big).
  \label{eq:rho-regret}
\end{equation}
Spearman rank correlation \cite{spearman1904} is preferred to a parametric correlation because the
relationship between a validation distance and a bounded success rate is monotone but not linear.
$\rho_m$ near $+1$ indicates that the metric ranks checkpoints in the same order as closed-loop success;
$\mathrm{Regret}(m)$ is the success forgone, in absolute percentage points, by trusting the metric's
selection $k_m^\star$ instead of the checkpoint that was in fact best. A good metric has high $\rho_m$
and low regret; the two need not agree in every instance, since regret depends on where $Q$ happens to be
noisy near $k_m^\star$, while $\rho_m$ aggregates over the whole curve.

\subsection{The Plateau Regime: A Restriction-of-Range Diagnostic}
\label{sec:method-plateau}

A single $\rho_m$ computed over an entire training run conflates two qualitatively different regimes: an
early regime in which both $\bar m(k)$ and $Q(k)$ change rapidly as the policy leaves initialization, and
a later, converged regime in which $\bar m(k)$ has plateaued and any residual variation in $Q(k)$ is
plausibly dominated by simulator evaluation noise. To separate the two, each run's \emph{knee} is defined
as the first checkpoint at which DTW has completed $90\%$ of its total descent,
\begin{equation}
  k_{\mathrm{knee}} = \min\Big\{k : \bar m_{\mathrm{DTW}}(k) \le m_{\min} + 0.1\big(\bar
  m_{\mathrm{DTW}}(0)-m_{\min}\big)\Big\}, \qquad m_{\min}=\min_j \bar m_{\mathrm{DTW}}(j).
  \label{eq:knee}
\end{equation}
$\rho_{\mathrm{DTW}}$ is then recomputed twice: once over the full evaluated range, and once restricted
to $\{k\ge k_{\mathrm{knee}}\}$, the \emph{converged plateau}. A metric that is only a coarse
regime-separator (untrained vs. converged) will show a large drop between the two; a metric that retains
fine-grained discriminative power inside the plateau will not. This diagnostic is what allows the
results in \S\ref{sec:results-plateau} to distinguish ``the metric stopped tracking success'' from ``the
evaluation schedule under-sampled a noisy, low-signal region.''

\subsection{Statistical Inference: Bootstrap Confidence Intervals}
\label{sec:method-bootstrap}

Every $\rho_m$ above is a point estimate computed from as few as $11$ and at most $33$ simulator
evaluations per run, each itself a $\pm5\%$-noisy $100$-episode Bernoulli success-rate estimate. Reading
an ordering directly off point estimates (``metric $A$ beats metric $B$'') risks over-interpreting
sampling noise. We attach a percentile bootstrap confidence interval to every $\rho_m$
\cite{efron1979bootstrap,efron1993bootstrap}: for each run, its evaluated checkpoints are resampled with
replacement $B=10{,}000$ times, a tie-corrected Spearman correlation is recomputed on each resample, and
the central $95\%$ interval is reported together with $\Pr(\rho_m>0)$. To compare two metrics within a
run without their mutual correlation confounding the test, the \emph{same} bootstrap resample is used for
both and the paired difference $\Delta\rho=\rho_A-\rho_B$ is reported over the checkpoints they share.
This yields three tiers of claim, from weakest to strongest: a metric's point-estimate $\rho_m$; whether
that $\rho_m$'s CI excludes zero (a significant association with success); and whether a paired
$\Delta\rho$ CI excludes zero (a significant \emph{ordering} between two metrics). Only the last licenses
a claim of the form ``metric $A$ is better than metric $B$.''

\section{Experiment Implementation}
\label{sec:implementation}

\subsection{Source Protocol}
\label{sec:implementation-source}

The methodology of \S\ref{sec:method} was specified, prior to implementation, in two internal protocol
documents: an initial specification, \texttt{wenyan\_experiment\_instruction.tex}, and a subsequent
refinement, \texttt{refined\_experiment\_instruction.tex}, which formalizes the same protocol with
explicit dataset notation ($\mathcal D=\mathcal D_{\mathrm{train}}\cup\mathcal D_{\mathrm{val}}\cup
\mathcal D_{\mathrm{test}}$), two evaluation modes (Mode~A, action-chunk comparison, and Mode~B,
closed-loop state-action rollout against a learned world model), and two output-table templates (a
per-checkpoint \emph{metric curve table} and a per-metric \emph{ranking table}). Table~\ref{tab:impl-steps}
reproduces the refined protocol's six-step experimental procedure and states, for each step, how it was
realized in this study; the subsections that follow (\S\ref{sec:implementation-models}--\S\ref{sec:implementation-sim})
develop each realization in full.

\begin{table}[h]
\centering
\footnotesize
\caption{The six-step experimental procedure of the source protocol
(\texttt{refined\_experiment\_instruction.tex}, \S\ref{sec:implementation-source}) and its realization in
this study.}
\label{tab:impl-steps}
\begin{tabular}{p{1.6cm}p{5.6cm}p{6.3cm}}
\toprule
Step & Source protocol & Realized in this study \\
\midrule
1.\ Prepare data &
Split the data into training, validation, and test sets; compute action-normalization statistics from
the training set only; fix the evaluation horizon $H$ and window stride; save an identical set of
validation windows for use with every checkpoint. &
Stratified $80/10/10$ split by task under a fixed seed (\S\ref{sec:implementation-data}); $H{=}50$-step,
phase-aligned, non-overlapping windows shared across every checkpoint and metric; actions normalized per
dimension from training-split statistics only (\S\ref{sec:method-metrics}). \\
2.\ Fine-tune and save checkpoints &
Fine-tune each VLA model and save checkpoints $\theta_0,\theta_1,\ldots,\theta_K$ at fixed training
intervals. &
Five LoRA ($r{=}32$) fine-tuning runs across three architectures, logging $68$--$200$ checkpoints per run
at fixed step intervals (Table~\ref{tab:runs}, \S\ref{sec:implementation-runs}). \\
3.\ Action-chunk evaluation (Mode~A) &
For each checkpoint, generate action chunks on the validation windows --- teacher-forced pseudo-chunks
for one-step policies, native chunk outputs for chunking policies --- compute DTW, OT, cosine distance,
and NLL or its architecture-correct surrogate, and aggregate over validation windows. &
Implemented as specified for all three architectures (\S\ref{sec:implementation-mode}); DTW, OT, cosine,
and NLL logged under identical keys across models, supplemented with cross-entropy, token accuracy, and
$L_1$/$L_2$ where architecture-appropriate; aggregated at the trajectory level via the window-averaged
estimator selected in \S\ref{sec:results-taskdisc}. \\
4.\ World-model rollout evaluation (Mode~B) &
If a world model is available, validate it on held-out trajectories, roll out each checkpoint from
shared initial states, and compare the resulting state-action sequences with successful references using
DTW, OT, and cosine distance. &
\textbf{Not executed.} No world model was built or integrated; closed-loop rollout evaluation is
therefore outside this study's scope (\S\ref{sec:implementation-deviations}). \\
5.\ Choose checkpoints &
For each metric, compute the selected checkpoint $k_m^\star$ and compare the selections across metrics. &
$k_m^\star$ computed via Eq.~\ref{eq:selection} by \texttt{checkpoint\_selection.py} for every metric and
run (\S\ref{sec:implementation-pipeline}); selections compared in \S\ref{sec:results-permodel}
(Table~\ref{tab:rankings}). \\
6.\ Rank metrics &
Where downstream quality $Q(k)$ is available, rank metrics by Spearman correlation, selection regret,
and stability under bootstrap resampling. &
All three criteria computed for every run (\S\ref{sec:results-permodel}--\S\ref{sec:results-bootstrap});
bootstrap stability (\S\ref{sec:method-bootstrap}) extends the source protocol's single-line request into
a full paired-resampling analysis. \\
\bottomrule
\end{tabular}
\end{table}

Two components of this study fall outside the six-step procedure entirely, in the sense that neither
document requests them: the task-discriminative-validity prerequisite of \S\ref{sec:method-taskdisc},
established prior to Step~3 as a precondition for trusting any metric computed under it, and the
combined multi-metric score of the source protocol's optional extension, which is not computed here
since every metric in \S\ref{sec:results} is ranked individually.

\subsection{Models}
\label{sec:implementation-models}

Three VLA architectures were fine-tuned, spanning both major action-generation paradigms in the current
literature. \textbf{OpenVLA} \cite{kim2024openvla} is a $7$B-parameter discrete-token policy built on a
Llama-2 backbone with a fused DINOv2/SigLIP visual encoder, autoregressively decoding actions as discrete
tokens. \textbf{SmolVLA} \cite{shukor2025smolvla} is a $450$M-parameter flow-matching policy designed for
efficient training and inference on consumer hardware. \textbf{$\pi_{0.5}$} \cite{black2025pi05} is a
$3$B-parameter flow-matching policy from Physical Intelligence trained for open-world generalization. All
three were adapted with Low-Rank Adaptation \cite{hu2021lora} at rank $r{=}32$, so that fine-tuning
touches a small, comparable parameter budget across architectures of very different scale.

\subsection{Benchmark and Task Suite}
\label{sec:implementation-benchmark}

All experiments use the LIBERO-goal task suite of the LIBERO benchmark \cite{liu2023libero}, ten
manipulation tasks sharing a common scene and object set but differing in goal condition, in the
MuJoCo-based LIBERO simulator. Closed-loop success $Q(k)$ is the fraction of successful rollouts over
$100$ simulated episodes per checkpoint, spread across the ten tasks; at this sample size a single $Q(k)$
estimate carries roughly $\pm5\%$ sampling noise, a figure used throughout \S\ref{sec:results} to
interpret apparent differences between checkpoints.

\subsection{Data Splits and Normalization}
\label{sec:implementation-data}

Each task's demonstrations are split $80/10/10$ into train/validation/test, stratified by task and fixed
by a common random seed so that the same validation and test trajectories are held out across every
run reported here, including the two runs deliberately trained on a $150$-trajectory subset of the
training split (Table~\ref{tab:runs} below): capping only the \emph{training} split at $150$
trajectories leaves the $44$ validation and $42$ test trajectories (on the $342$-trajectory task set)
unchanged, so that validation-set metrics remain comparable between the $150$-trajectory and full-data
runs. Actions are normalized per dimension using statistics computed on the training split only
(\S\ref{sec:method-metrics}).

\subsection{Fine-Tuning Configuration and the Five Runs}
\label{sec:implementation-runs}

Five fine-tuning runs are reported, summarized in Table~\ref{tab:runs}. The first four form a
$2\times2$ design crossing architecture (discrete-token OpenVLA vs.\ flow-matching SmolVLA/$\pi_{0.5}$)
with training-set size ($150$ vs.\ the full $342$ trajectories); the fifth, $\pi_{0.5}$-150, is a
deliberate overtraining \emph{stress} run using an aggressive learning rate
($3\times10^{-4}$, roughly $6\times$ the other runs') on the smaller $150$-trajectory subset, trained
substantially past where the other runs stop.

\begin{table}[h]
\centering
\caption{The five fine-tuning runs. All LIBERO-goal, LoRA $r{=}32$; $Q(k)$ is $100$-episode simulated
success.}
\label{tab:runs}
\begin{tabular}{lllrrr}
\toprule
Run & Architecture & Train traj. & Checkpoints & Sim.\ evals & Peak $Q$ \\
\midrule
SmolVLA-150      & flow-matching (450M) & 150 & 200 (to 100k) & 18 & 0.68 \\
OpenVLA-150      & discrete-token (7B)  & 150 & 79 (to 39.5k) & 19 & 0.83 \\
OpenVLA-full     & discrete-token (7B)  & 342 & 68 (to 34k)   & 11 & 0.78 \\
$\pi_{0.5}$-full & flow-matching (3B)   & 342 & 148 (to 73.5k)& 32 & 0.18$^\dagger$ \\
$\pi_{0.5}$-150  & flow-matching (3B)   & 150 & 80 (to 40k)   & 28 & 0.38$^\dagger$ \\
\bottomrule
\end{tabular}
\\[2pt]
{\footnotesize $^\dagger$ Both $\pi_{0.5}$ runs' success is measured at the confounded open-loop action
horizon $n_{\mathrm{action}}{=}50$ (the full predicted chunk is executed blind before replanning), which
roughly halves LIBERO success relative to shorter replanning horizons; only the \emph{relative} ranking
of metrics is meaningful for these two rows, not the absolute success level.}
\end{table}

\subsection{Evaluation Mode and Metric Instantiation}
\label{sec:implementation-mode}

Consistent with Mode~A of the source protocol (\S\ref{sec:implementation-source}), OpenVLA's action
chunks are teacher-forced (AC-TF) and SmolVLA's and $\pi_{0.5}$'s are native chunk predictions
(AC-Chunk). Every $\bar m(k)$ reported in \S\ref{sec:results} uses the trajectory-level, window-averaged
estimator (B) of \S\ref{sec:method-taskdisc}, selected over the single-chunk estimator (A) on the basis
of the task-discriminative-validity comparison reported in \S\ref{sec:results-taskdisc}.

\subsection{Offline Metric Computation Pipeline}
\label{sec:implementation-pipeline}

Offline metrics are computed and logged during fine-tuning at fixed validation intervals by the
fine-tuning scripts themselves, and cached per run in a \texttt{metrics\_history.json} record of
$\{\text{step}, \{\text{metric}:\text{value}\}\}$. A dedicated \texttt{checkpoint\_selection.py} module
(kept identical, i.e.\ ``in lockstep,'' between the OpenVLA and the SmolVLA/$\pi_{0.5}$ code paths) then
implements Eq.~\ref{eq:selection} exactly, and a companion \texttt{select\_and\_eval\_checkpoints.py}
script shells out to each framework's own simulator-evaluation entry point --- \texttt{run\_libero\_eval.py}
for OpenVLA, the \texttt{lerobot-eval} CLI for the LoRA/PEFT SmolVLA and $\pi_{0.5}$ checkpoints --- to
obtain $Q(k_m^\star)$ for the checkpoint each metric selects, caching results so that repeated analysis
does not require repeated simulation.

\subsection{Simulator Evaluation Protocol}
\label{sec:implementation-sim}

Because a full sweep of $Q(k)$ over every logged checkpoint is prohibitively expensive, simulator
evaluation is sparse by design: only a subset of checkpoints per run (Table~\ref{tab:runs}, ``Sim.\
evals'') are ever rolled out in simulation, each for $100$ episodes across the ten LIBERO-goal tasks.
This sparsity is a first-order constraint on everything reported in \S\ref{sec:results}: it is the reason
$\rho_m$ must be treated as an uncertain statistic (\S\ref{sec:method-bootstrap}) rather than a fact, and
the reason the plateau diagnostic of \S\ref{sec:method-plateau} exists at all.

\subsection{Deviations from and Extensions to the Source Protocol}
\label{sec:implementation-deviations}

Table~\ref{tab:impl-steps} realizes five of the source protocol's six steps; Step~4 is the one
structural gap and warrants comment. No world model was built or integrated, so Mode~B (closed-loop
rollout evaluation) is entirely out of scope: every result in this chapter is an open-loop,
teacher-forced-or-native-chunk comparison, and the qualitative claim that ``open-loop action similarity
is not the same thing as closed-loop rollout quality'' is asserted from the literature rather than
measured directly in this study. Relatedly, the source protocol's optional combined multi-metric score
is not computed; every metric in \S\ref{sec:results} is ranked individually, and no attempt is made to
fuse metrics into a single decision rule.

\section{Experimental Results}
\label{sec:results}

\subsection{Task-Discriminative Validity}
\label{sec:results-taskdisc}

Before any metric's correlation with success is evaluated, its ability to separate same-task from
different-task motion (\S\ref{sec:method-taskdisc}) is checked on $428$ held-out ground-truth
trajectories spanning the ten LIBERO-goal tasks, using $3{,}000$ intra-task and $3{,}000$ inter-task
chunk pairs. Table~\ref{tab:taskdisc} reports both estimators' results side by side.

\begin{table}[h]
\centering
\caption{Task-discriminative validity: single-chunk estimator (A) vs.\ trajectory-level, window-averaged
estimator (B), $3{,}000$ cross-trajectory pairs each. Sorted by estimator-(B) AUROC.}
\label{tab:taskdisc}
\begin{tabular}{lrrrr}
\toprule
 & \multicolumn{2}{c}{estimator (A): single chunk} & \multicolumn{2}{c}{estimator (B): trajectory-level} \\
\cmidrule(lr){2-3}\cmidrule(lr){4-5}
Metric & AUROC & Cohen's $d$ & AUROC & Cohen's $d$ \\
\midrule
Cosine    & 0.657 & 0.66 & 0.927 & 2.19 \\
OT        & 0.692 & 0.62 & 0.923 & 1.68 \\
DTW       & 0.677 & 0.47 & 0.911 & 1.61 \\
$L_2$     & 0.653 & 0.39 & 0.903 & 1.70 \\
$L_1$     & 0.660 & 0.56 & 0.901 & 1.80 \\
\bottomrule
\end{tabular}
\end{table}

Every metric clears chance under both estimators, but the two differ sharply in strength. The single
chunk (A) separates tasks only moderately (AUROC $0.65$--$0.69$), because LIBERO-goal's ten tasks share a
large common motion vocabulary --- reach, grasp, transport --- so an arbitrary $50$-step window overlaps
heavily across tasks. The trajectory-level, phase-aligned estimator (B) is decisively stronger: every
metric rises to AUROC $\ge0.90$ and Cohen's $d$ roughly triples. The improvement is attributable
specifically to phase alignment rather than variance reduction alone --- e.g.\ DTW's intra-task mean
distance falls from $7.06$ to $1.84$ under (B) while its inter-task mean barely moves ($10.23\to8.27$),
consistent with comparing corresponding motion phases rather than arbitrary ones. On this basis,
estimator (B) is adopted for every $\bar m(k)$ in the remainder of this chapter; the precondition of
\S\ref{sec:method-taskdisc} is satisfied for all five candidate metrics.

\subsection{Cross-Model Correlation with Closed-Loop Success}
\label{sec:results-crossmodel}

\begin{figure}[h]
  \centering
  \includegraphics[width=0.92\textwidth]{figures/crossmodel_early_stopping.png}
  \caption{Spearman $\rho_m$ of every offline metric with LIBERO-goal success, by run (left), and the
  corresponding success curve $Q(k)$ over training (right). Trajectory (DTW/OT/cosine) and regression
  ($L_1$/$L_2$) metrics are positive in every run; the likelihood/loss family (NLL, or the flow-matching
  loss standing in for it) splits by architecture, positive and significant for the two discrete-token
  runs, weak or statistically indistinguishable from zero for the flow-matching runs
  (\S\ref{sec:results-bootstrap}).}
  \label{fig:crossmodel}
\end{figure}

Figure~\ref{fig:crossmodel} summarizes the headline result across all five runs and all applicable
metrics. Ordering the five runs SmolVLA-150 / OpenVLA-150 / OpenVLA-full / $\pi_{0.5}$-full /
$\pi_{0.5}$-150, the trajectory-similarity metrics are positive throughout: $\rho_{\mathrm{DTW}} =
0.47/0.75/0.85/0.90/0.72$, $\rho_{\mathrm{OT}} = 0.34/0.84/0.66/0.89/0.68$, and $\rho_{\cos} =
0.28/0.75/0.89/0.86/0.77$. DTW in particular is never negative across the five runs and its selection
regret never exceeds $0.17$, a first empirical basis for the conclusion drawn in
\S\ref{sec:results-conclusions}\ref{conc:correlation}.

\subsection{Per-Model Metric Rankings}
\label{sec:results-permodel}

Table~\ref{tab:rankings} reports, per run, every applicable metric's selected checkpoint $k_m^\star$,
Spearman correlation $\rho_m$, and selection regret, sorted by $\rho_m$.

\begin{table}[h]
\centering
\caption{Per-run metric rankings: selected checkpoint, Spearman $\rho_m$, and selection regret
(Eq.~\ref{eq:selection}--\ref{eq:rho-regret}).}
\label{tab:rankings}
\begin{minipage}[t]{0.48\textwidth}
\centering \textbf{SmolVLA-150} (flow-matching)\\[2pt]
\begin{tabular}{lrrr}
\toprule
Metric & $k_m^\star$ & $\rho$ & Regret \\ \midrule
DTW    & 22000 & $+0.47$ & 0.02 \\
$L_2$  & 24500 & $+0.46$ & 0.05 \\
$L_1$  & 56000 & $+0.44$ & 0.24 \\
OT     & 11000 & $+0.34$ & 0.26 \\
Cosine & 33000 & $+0.28$ & 0.08 \\
NLL/FM & 13000 & $-0.22$ & 0.19 \\
\bottomrule
\end{tabular}
\end{minipage}
\hfill
\begin{minipage}[t]{0.48\textwidth}
\centering \textbf{OpenVLA-150} (discrete-token)\\[2pt]
\begin{tabular}{lrrr}
\toprule
Metric & $k_m^\star$ & $\rho$ & Regret \\ \midrule
OT     & 33499 & $+0.84$ & 0.08 \\
Acc.   & 37499 & $+0.78$ & 0.02 \\
$L_1$  & 30999 & $+0.78$ & 0.06 \\
DTW    & 23499 & $+0.75$ & 0.17 \\
Cosine & 23499 & $+0.75$ & 0.17 \\
$L_2$  & 23499 & $+0.75$ & 0.17 \\
CE     & 30999 & $+0.72$ & 0.06 \\
NLL    & 30999 & $+0.72$ & 0.06 \\
\bottomrule
\end{tabular}
\end{minipage}

\vspace{1em}

\begin{minipage}[t]{0.48\textwidth}
\centering \textbf{OpenVLA-full} (342 traj.)\\[2pt]
\begin{tabular}{lrrr}
\toprule
Metric & $k_m^\star$ & $\rho$ & Regret \\ \midrule
Cosine & 28499 & $+0.89$ & 0.05 \\
DTW    & 21499 & $+0.85$ & 0.00 \\
$L_2$  & 21499 & $+0.85$ & 0.00 \\
$L_1$  & 29999 & $+0.85$ & 0.05 \\
CE     & 33999 & $+0.83$ & 0.04 \\
Acc.   & 33999 & $+0.83$ & 0.04 \\
NLL    & 33999 & $+0.83$ & 0.04 \\
OT     & 33999 & $+0.66$ & 0.04 \\
\bottomrule
\end{tabular}
\end{minipage}
\hfill
\begin{minipage}[t]{0.48\textwidth}
\centering \textbf{$\pi_{0.5}$-full} (flow-matching, 342 traj.)\\[2pt]
\begin{tabular}{lrrr}
\toprule
Metric & $k_m^\star$ & $\rho$ & Regret \\ \midrule
DTW    & 57999 & $+0.90$ & 0.04 \\
OT     & 60499 & $+0.89$ & 0.04 \\
Cosine & 55999 & $+0.86$ & 0.04 \\
$L_1$  & 55999 & $+0.84$ & 0.04 \\
$L_2$  & 55999 & $+0.81$ & 0.04 \\
NLL/FM & 31999 & $+0.46$ & 0.08 \\
\bottomrule
\end{tabular}
\end{minipage}

\vspace{1em}
\centering
\begin{minipage}[t]{0.48\textwidth}
\centering \textbf{$\pi_{0.5}$-150} (overtraining stress run)\\[2pt]
\begin{tabular}{lrrr}
\toprule
Metric & $k_m^\star$ & $\rho$ & Regret \\ \midrule
$L_1$  & 31999 & $+0.78$ & 0.05 \\
Cosine & 39499 & $+0.77$ & 0.04 \\
$L_2$  & 37999 & $+0.73$ & 0.08 \\
DTW    & 33499 & $+0.72$ & 0.07 \\
OT     & 24499 & $+0.68$ & 0.17 \\
NLL/FM &  8999 & $-0.16$ & 0.20 \\
\bottomrule
\end{tabular}
\end{minipage}
\end{table}

Two mechanistic points follow from Table~\ref{tab:rankings}. First, the metrics available depend on
architecture: only OpenVLA exposes cross-entropy loss and token accuracy, since only a discrete-token
policy has them; the flow-matching policies retain $L_1$/$L_2$ regression error, which is on par with or
above their own DTW score in every flow-matching run. Second, on OpenVLA-full, $L_2$ and DTW select the
\emph{identical} checkpoint ($k^\star=21499$) with identical $\rho$ and zero regret, foreshadowing the
statistical indistinguishability of the trajectory and regression families confirmed in
\S\ref{sec:results-bootstrap}.

\subsection{The Plateau Regime and Restriction of Range}
\label{sec:results-plateau}

\begin{table}[h]
\centering
\caption{DTW$\to$success Spearman $\rho$ over the full evaluated range vs.\ restricted to the converged
plateau $\{k\ge k_{\mathrm{knee}}\}$ (Eq.~\ref{eq:knee}).}
\label{tab:plateau}
\begin{tabular}{lrrrrr}
\toprule
Run & $k_{\mathrm{knee}}$ & Plateau pts & $\rho_{\mathrm{full}}$ & $\rho_{\mathrm{plateau}}$ & Drop \\
\midrule
SmolVLA-150      &  4500 & 17/18 & $+0.47$ & $+0.38$ & 0.10 \\
OpenVLA-150      &  6999 & 16/19 & $+0.75$ & $+0.59$ & 0.16 \\
OpenVLA-full     &  8499 &  9/11 & $+0.86$ & $+0.73$ & 0.12 \\
$\pi_{0.5}$-full & 18999 & 22/32 & $+0.91$ & $+0.80$ & 0.11 \\
$\pi_{0.5}$-150  &  9999 & 22/28 & $+0.72$ & $\mathbf{+0.43}$ & $\mathbf{0.29}$ \\
\bottomrule
\end{tabular}
\end{table}

\begin{figure}[h]
  \centering
  \includegraphics[width=0.92\textwidth]{figures/plateau_restriction.png}
  \caption{\textbf{Left:} DTW$\to$success $\rho$, full range vs.\ converged plateau, all five runs.
  \textbf{Right:} $\pi_{0.5}$-150's validation DTW (converges early, axis inverted) against its success
  curve (a noisy, low-ceiling band that dips near $20$k training steps and recovers to a $38\%$ peak by
  $34$k once training is extended to $40$k steps).}
  \label{fig:plateau}
\end{figure}

$\rho_{\mathrm{full}}\ge\rho_{\mathrm{plateau}}$ holds for every run (Table~\ref{tab:plateau}): the
offline metrics are, at minimum, reliable regime separators (untrained vs.\ converged), and fine-ranking
already-converged checkpoints is strictly harder. For four of the five runs the plateau correlation
remains healthy ($\rho_{\mathrm{plateau}}\ge0.59$, drop $\le0.16$), so the metric retains real
discriminative power even inside the plateau. $\pi_{0.5}$-150 is the exception, dropping to
$\rho_{\mathrm{plateau}}=+0.43$. Critically, this is a restriction-of-range artifact rather than genuine
action-quality overtraining: an earlier $20{,}000$-step view of this run showed success falling from
$31\%$ to $23\%$ after its peak, consistent with overtraining, but extending training to $40{,}000$
steps shows success \emph{recovering} to a $38\%$ peak (Figure~\ref{fig:plateau}, right) --- so the
apparent decline was an artifact of stopping mid-dip in a low-ceiling, noisy success band whose spread is
comparable to the $\pm5\%$ sampling noise of a $100$-episode evaluation, not a monotone downturn. The
practical implication, developed into a conclusion in \S\ref{sec:results-conclusions}\ref{conc:plateau},
is that these metrics support a coarse but robust early-stopping decision (do not stop before the knee)
even where they cannot reliably fine-rank checkpoints once training has converged.

\subsection{Statistical Power: Bootstrap Confidence Intervals}
\label{sec:results-bootstrap}

\begin{table}[h]
\centering
\caption{Bootstrap $95\%$ confidence intervals on $\rho$ ($B{=}10{,}000$ resamples, tie-corrected
Spearman, paired checkpoint resampling per Eq.~\ref{eq:rho-regret}\S\ref{sec:method-bootstrap}). ``\#sig''
counts metrics whose interval excludes $0$; ``best $L$'' is the stronger of $L_1$/$L_2$.}
\label{tab:bootstrap}
\begin{tabular}{lccccc}
\toprule
Run & $n$ & \#sig & DTW $\rho$ [95\% CI] & best $L$ $\rho$ [95\% CI] & NLL $\rho$ [95\% CI] \\
\midrule
SmolVLA-150      & 18 & 0/6 & $0.47\,[-0.09,\,0.87]$ & $0.45\,[-0.12,\,0.84]$ & $-0.21\,[-0.74,\,0.35]$ \\
OpenVLA-150      & 19 & 8/8 & $0.75\,[0.34,\,0.94]$  & $0.78\,[0.42,\,0.95]$  & $0.72\,[0.33,\,0.92]$ \\
OpenVLA-full     & 11 & 8/8 & $0.84\,[0.46,\,0.99]$  & $0.84\,[0.42,\,0.98]$  & $0.81\,[0.40,\,0.97]$ \\
$\pi_{0.5}$-full & 33 & 6/6 & $0.90\,[0.79,\,0.95]$  & $0.83\,[0.64,\,0.93]$  & $0.41\,[0.01,\,0.70]$ \\
$\pi_{0.5}$-150  & 28 & 5/6 & $0.71\,[0.37,\,0.90]$  & $0.78\,[0.50,\,0.93]$  & $-0.17\,[-0.57,\,0.25]$ \\
\bottomrule
\end{tabular}
\end{table}

Statistical power, not the metric family, is the binding constraint on how strong a claim the data
support. Only $\pi_{0.5}$-full ($n{=}33$) is powered enough to resolve within-run structure; at the
other extreme, no metric on SmolVLA-150 ($n{=}18$) clears zero at $95\%$, including DTW itself
($0.47$, CI $[-0.09,0.87]$) --- consistent with \S\ref{sec:results-plateau}'s finding that this run's
evaluation schedule sampled almost only the post-knee plateau. Within every run, DTW's interval overlaps
OT's, cosine's, $L_1$'s, and $L_2$'s: the paired-resampling test finds DTW rank-\emph{identical} to
$L_2$ on OpenVLA-full ($\Delta\rho\equiv0$) and finds $\Pr(\rho_{\mathrm{DTW}}>\rho_{L_2})$
indistinguishable from a coin flip on the three intermediate runs ($0.63/0.44/0.33$). DTW significantly
outperforms a loss-based metric in exactly one comparison: against the flow-matching surrogate NLL/FM on
$\pi_{0.5}$-full ($\Pr=1.00$, $\Delta\rho$ CI $[0.26,0.92]$). The architecture split in NLL's validity is
statistically confirmed for OpenVLA (significantly positive in both OpenVLA runs) and, more weakly, for
$\pi_{0.5}$-full; for SmolVLA-150 and $\pi_{0.5}$-150 the NLL/FM interval spans zero, so the negative
point estimates in those two rows should be read as \emph{uninformative}, not as evidence of a reliably
anti-correlated surrogate.

\subsection{Conclusions}
\label{sec:results-conclusions}

The results of \S\ref{sec:results-taskdisc}--\S\ref{sec:results-bootstrap} support the following
conclusions.

\begin{enumerate}[label=(\alph*)]
  \item \textbf{Offline trajectory-similarity metrics correlate strongly with closed-loop success.}
    \label{conc:correlation}
    DTW, OT, and cosine distance are positive in every one of the five runs, spanning both
    architectures (discrete-token and flow-matching) and both training-set sizes
    (Figure~\ref{fig:crossmodel}, Table~\ref{tab:rankings}). DTW's association is significant (bootstrap
    $95\%$ CI excluding zero) in four of five runs and is never negative on point estimates, making it
    the most consistently reliable default success predictor in this study, though
    \S\ref{sec:results-bootstrap} shows it is not statistically separable from the regression losses
    $L_1$/$L_2$ or from OT/cosine.

  \item \textbf{The metrics reliably identify the converged (plateau) region and support low-regret
    early-stop checkpoint selection, though fine-ranking inside an already-converged plateau is
    harder.} \label{conc:plateau}
    Every metric's selection regret stays $\le0.26$ across all five runs and $\le0.17$ for DTW
    specifically (Table~\ref{tab:rankings}), and the coarse decision ``do not stop before the knee'' is
    supported by $\rho\approx0.8$--$0.9$ over the full descent in every run. Restricted to the converged
    plateau, $\rho$ weakens for every run (Table~\ref{tab:plateau}) but remains usable
    ($\rho_{\mathrm{plateau}}\ge0.59$) for four of five; only $\pi_{0.5}$-150 loses most of its
    plateau-region signal, and this is attributable to a noisy, low success ceiling rather than to a
    metric failure (\S\ref{sec:results-plateau}).

  \item \textbf{Likelihood- and loss-based surrogates are reliable only for discrete-token
    architectures.} \label{conc:nll}
    OpenVLA's NLL is a genuine token log-likelihood and a strong, significant predictor in both of its
    runs. For the flow-matching policies, the same-named surrogate is the validation flow-matching loss,
    and it is uniformly weaker, significant only on $\pi_{0.5}$-full, and statistically indistinguishable
    from zero on the two smaller flow-matching runs (\S\ref{sec:results-bootstrap}). A loss-based early
    stop is therefore defensible for discrete-token policies but not recommended for flow-matching
    policies, for which a trajectory-similarity metric should be used instead.

  \item \textbf{Discrete-token architectures expose additional usable offline signals, without any
    single metric dominating.} \label{conc:signals}
    OpenVLA additionally exposes cross-entropy loss and token accuracy, both strong predictors on both
    training-set sizes; on the full-data run, $L_2$ is rank-identical to DTW at zero regret. The paired
    bootstrap test confirms this is not a point-estimate coincidence: the regression losses are
    statistically indistinguishable from DTW in every run, so the applicable metric set --- DTW, OT,
    cosine, $L_1$, $L_2$ --- should be read as a family of comparably valid choices rather than a strict
    ranking.

  \item \textbf{No catastrophic action-quality overtraining is observed in any run, including the two
    stress runs.} \label{conc:overtrain}
    Both $150$-trajectory runs (OpenVLA and SmolVLA) were designed to induce overtraining relative to
    their full-data counterparts; instead, every run settles into a noisy plateau in both the offline
    metrics and $Q(k)$ with no monotone late decline. $\pi_{0.5}$-150, trained with an aggressive
    learning rate specifically to stress-test this claim, initially appeared to be the exception at a
    $20$k-step horizon, but extending training to $40$k steps shows the apparent decline was a
    restriction-of-range artifact, not genuine overtraining (\S\ref{sec:results-plateau}). In-distribution
    validation together with LoRA's limited parameter budget appears sufficient, in this setting, to
    suppress action-space overfitting even under reduced data and an aggressive learning rate.

  \item \textbf{At the sample sizes available here, statistical power --- not the identity of the
    metric --- is usually the binding constraint on what can be claimed.} \label{conc:power}
    With $11$--$33$ simulator evaluations per run, most within-run \emph{orderings} among the
    trajectory and regression metrics are not statistically resolved (\S\ref{sec:results-bootstrap});
    only the architecture split of conclusion~\ref{conc:nll} and the positive-vs.-uninformative
    distinction of the likelihood surrogate survive the confidence intervals. Conclusions~\ref{conc:correlation}--\ref{conc:overtrain}
    should therefore be read at the level of metric \emph{families} (trajectory/regression vs.\
    likelihood-based) rather than as a strict ranking of individual metrics, and denser or lower-variance
    simulator evaluation would be required to resolve finer distinctions.
\end{enumerate}

\begin{equation}
  \boxed{\ \text{trajectory similarity (DTW/OT/cosine)} \;\succsim\; \text{regression loss}\;
  (L_1/L_2) \;\succ\; \text{flow-matching likelihood surrogate (NLL/FM)}\ }
  \label{eq:conclusion-box}
\end{equation}

\section{Summary}
\label{sec:summary}

\paragraph{Motivation and research question.} Robot foundation models are expensive to fine-tune and
expensive to evaluate online, yet training loss alone does not reliably indicate when a checkpoint has
reached its best closed-loop performance (\S\ref{sec:intro}). This chapter asked which \emph{offline}
validation metric --- computed cheaply from held-out demonstrations at every checkpoint, without any
robot or simulator rollout --- best predicts the closed-loop success rate that would otherwise require
one.

\paragraph{Method and implementation.} \S\ref{sec:method} formalized four candidate metrics (DTW, OT,
cosine distance, and an architecture-appropriate likelihood or loss surrogate), established
task-discriminative validity as a precondition for trusting any of them, and defined the early-stopping
rule, rank-correlation/regret scoring, a plateau/knee diagnostic for restriction-of-range effects, and a
bootstrap procedure for quantifying statistical confidence. \S\ref{sec:implementation} instantiated this
protocol --- following, and where noted extending, the \texttt{wenyan}/\texttt{refined}
\texttt{\_experiment\_instruction.tex} source protocol (Table~\ref{tab:impl-steps}) --- across five LoRA
fine-tuning runs spanning three architectures (discrete-token OpenVLA; flow-matching SmolVLA and
$\pi_{0.5}$), two training-set sizes, and one deliberate overtraining stress run, evaluated on the
LIBERO-goal benchmark.

\paragraph{Results.} \S\ref{sec:results} found that trajectory-similarity metrics --- DTW above all ---
correlate positively with closed-loop success in every run, regardless of architecture or training-set
size, and are not statistically distinguishable from the regression losses $L_1$/$L_2$ at the sample
sizes available (conclusions~\ref{conc:correlation}, \ref{conc:signals}). These metrics support a coarse
but robust early-stopping decision --- avoid stopping before the descent knee --- and, for four of the
five runs, retain usable discriminative power even inside the converged plateau
(conclusion~\ref{conc:plateau}). The one run where plateau-region signal collapses, $\pi_{0.5}$-150, does
so because of a noisy, low success ceiling rather than because the metric fails, a distinction only the
plateau diagnostic and the extended-horizon re-evaluation of \S\ref{sec:results-plateau} could establish.
Likelihood- and loss-based surrogates track success only for the discrete-token architecture
(conclusion~\ref{conc:nll}); none of the five runs, including the two designed to stress-test it, shows
catastrophic action-quality overtraining (conclusion~\ref{conc:overtrain}); and bootstrap resampling
showed that statistical power, not the identity of the metric, is usually what limits how finely two
metrics can be ordered against each other (conclusion~\ref{conc:power}).

\paragraph{Limitations and directions for future chapters.} Three scope limitations, already flagged
where they arise, bound how far these conclusions generalize. First, Mode~B of the source protocol ---
closed-loop rollout evaluation against a learned world model --- was never implemented
(\S\ref{sec:implementation-deviations}), so every result in this chapter is an open-loop comparison; a
world-model-based closed-loop study is the most direct way to test whether the metrics ranked here also
predict compounding rollout error, not only single-chunk action quality. Second, every metric was
evaluated on the validation split $\mathcal D_{\mathrm{val}}$ used to select checkpoints; a held-out
$\mathcal D_{\mathrm{test}}$ split already exists in the data pipeline (\S\ref{sec:implementation-data})
but was not used to confirm these rankings replicate out of sample. Third, closed-loop success itself was
measured sparsely (11--33 simulator evaluations per run) and, for both $\pi_{0.5}$ runs, at a confounded
action-execution horizon (Table~\ref{tab:runs}); denser, lower-variance evaluation at a matched horizon
across architectures would sharpen the bootstrap intervals of \S\ref{sec:results-bootstrap} and could
resolve orderings --- such as DTW vs.\ $L_1$/$L_2$ --- that the present sample sizes leave statistically
open. None of these limitations changes the chapter's central claim: an offline trajectory-similarity
metric, computed once per checkpoint from held-out demonstrations, is a practical and
architecture-agnostic signal for deciding when to stop fine-tuning a robot foundation model.

\begin{thebibliography}{99}

\bibitem{sakoe1978dtw}
H.~Sakoe and S.~Chiba, ``Dynamic programming algorithm optimization for spoken word recognition,''
\emph{IEEE Transactions on Acoustics, Speech, and Signal Processing}, vol.~26, no.~1, pp.~43--49, 1978.

\bibitem{cuturi2013sinkhorn}
M.~Cuturi, ``Sinkhorn distances: Lightspeed computation of optimal transport,'' in \emph{Advances in
Neural Information Processing Systems (NeurIPS)}, 2013.

\bibitem{spearman1904}
C.~Spearman, ``The proof and measurement of association between two things,'' \emph{American Journal of
Psychology}, vol.~15, no.~1, pp.~72--101, 1904.

\bibitem{efron1979bootstrap}
B.~Efron, ``Bootstrap methods: Another look at the jackknife,'' \emph{The Annals of Statistics}, vol.~7,
no.~1, pp.~1--26, 1979.

\bibitem{efron1993bootstrap}
B.~Efron and R.~J.~Tibshirani, \emph{An Introduction to the Bootstrap}. Chapman \& Hall/CRC, 1993.

\bibitem{goodfellow2016deep}
I.~Goodfellow, Y.~Bengio, and A.~Courville, \emph{Deep Learning}. MIT Press, 2016.

\bibitem{ying2019overview}
X.~Ying, ``An overview of overfitting and its solutions,'' \emph{Journal of Physics: Conference Series},
vol.~1168, article 022022, 2019.

\bibitem{hu2021lora}
E.~J.~Hu, Y.~Shen, P.~Wallis, Z.~Allen-Zhu, Y.~Li, S.~Wang, L.~Wang, and W.~Chen, ``LoRA: Low-rank
adaptation of large language models,'' \emph{arXiv preprint arXiv:2106.09685}, 2021.

\bibitem{vlasurvey}
M.~Ud~Din, W.~Akram, L.~Saad~Saoud, J.~Rosell, and I.~Hussain, ``Vision language action models in robotic
manipulation: A systematic review,'' \emph{arXiv preprint arXiv:2507.10672}, 2025.

\bibitem{kim2024openvla}
M.~J.~Kim, K.~Pertsch, S.~Karamcheti, T.~Xiao, A.~Balakrishna, S.~Nair, R.~Rafailov, E.~Foster, G.~Lam,
P.~Sanketi, Q.~Vuong, T.~Kollar, B.~Burchfiel, R.~Tedrake, D.~Sadigh, S.~Levine, P.~Liang, and C.~Finn,
``OpenVLA: An open-source vision-language-action model,'' \emph{arXiv preprint arXiv:2406.09246}, 2024.

\bibitem{shukor2025smolvla}
M.~Shukor, D.~Aubakirova, F.~Capuano, et al., ``SmolVLA: A vision-language-action model for affordable
and efficient robotics,'' \emph{arXiv preprint arXiv:2506.01844}, 2025.

\bibitem{black2024pi0}
K.~Black, N.~Brown, D.~Driess, A.~Esmail, M.~Equi, C.~Finn, N.~Fusai, L.~Groom, K.~Hausman, B.~Ichter,
S.~Jakubczak, T.~Jones, L.~Ke, S.~Levine, A.~Li-Bell, M.~Mothukuri, S.~Nair, K.~Pertsch, L.~X.~Shi,
J.~Tanner, Q.~Vuong, A.~Walling, H.~Wang, and U.~Zhilinsky, ``$\pi_0$: A vision-language-action flow model
for general robot control,'' \emph{arXiv preprint arXiv:2410.24164}, 2024.

\bibitem{black2025pi05}
K.~Black, N.~Brown, D.~Driess, C.~Finn, S.~Levine, et al., ``$\pi_{0.5}$: A vision-language-action model
with open-world generalization,'' \emph{arXiv preprint arXiv:2504.16054}, 2025.

\bibitem{liu2023libero}
B.~Liu, Y.~Zhu, C.~Gao, Y.~Feng, Q.~Liu, Y.~Zhu, and P.~Stone, ``LIBERO: Benchmarking knowledge transfer
for lifelong robot learning,'' in \emph{Advances in Neural Information Processing Systems (NeurIPS),
Datasets and Benchmarks Track}, 2023. arXiv:2306.03310.

\bibitem{il}
A.~Hussein, M.~M.~Gaber, E.~Elyan, and C.~Jayne, ``Imitation learning: A survey of learning methods,''
\emph{ACM Computing Surveys}, vol.~50, no.~2, article 21, 2017.

\bibitem{dtw_il}
L.~Gong, B.~Chen, W.~Xu, C.~Liu, X.~Li, Z.~Zhao, and L.~Zhao, ``Motion similarity evaluation between
human and a Tri-Co robot during real-time imitation with a trajectory dynamic time warping model,''
\emph{Sensors}, vol.~22, no.~5, p.~1968, 2022.

\bibitem{ot_il}
R.~Dadashi, L.~Hussenot, M.~Geist, and O.~Pietquin, ``Primal Wasserstein imitation learning,'' in
\emph{International Conference on Learning Representations (ICLR)}, 2021. arXiv:2006.04678.

\bibitem{smoothness_il}
S.~Kulkarni, R.~Dhar, and Y.~Cui, ``Learning from the best: Smoothness-driven metrics for data quality in
imitation learning,'' \emph{arXiv preprint arXiv:2604.23000}, 2026.

\end{thebibliography}

\end{document}
